\section{Simulation methods}
\label{app:sim}

\paragraph{C1. Prices.} A Monte Carlo program in C (\texttt{sim\_v.c}, xoshiro256** generator, Box--Muller normals) advances all prices one block at a time with exact log-normal steps, $\ln p_{i,k+1}-\ln p_{i,k}=-\tfrac12\sigma_i^2\Delta t+\sigma_i\sqrt{\Delta t}\,(Lz_k)_i$, where $L$ is the Cholesky factor of the correlation matrix. Core-basket volatilities are 50, 65, 80, 100 and 80\%; correlations are BTC--ETH 0.80, BTC--SOL 0.70, BTC--HYPE 0.50, BTC--LINK 0.65, ETH--SOL 0.75, ETH--HYPE 0.50, ETH--LINK 0.75, SOL--HYPE 0.50, SOL--LINK 0.65 and HYPE--LINK 0.45. Runs cover one year with 12-second blocks, or 60 days (2\,s) and 12 days (0.4\,s), annualised. All designs in a scenario share the price path and the retail stream.

\paragraph{C2. Public pools.} In every block each risky token is arbitraged against USDC, in two passes over the tokens. The arbitrageur trades on the exact weighted-pool curve, with the fee charged on the input, until the pool's ask equals $p(1-c)$ or its bid equals $p(1+c)$, using the closed form from the pair invariant $x_i^{w_i}x_n^{w_n}$. The LP's loss is the change in pool value at true prices; turnover is the notional at true prices.

\paragraph{C3. Gated vaults.} At the end of each block the vault computes $D$ at true prices. If $D>\delta$ it sets weights to $w+(\din/D)(\omega-w)$ and deducts $s(D-\din)$ of NAV. In the price-selection stress the solver also chooses, for each leg, the least favourable price for the vault among the last $m$ blocks, and the vault loses that wedge on the leg's notional.

\paragraph{C4. Quote layer and informed trader.} Retail trades virtual reserves $\hat x_k=w_k\NAV/p_k$ with weighted-pool math and fee $f_q$; real balances change by the same amounts. The informed trader acts before retail in each block: for each token it trades toward the true price on a curve centred on an oracle price that lags by $\ell$ blocks, or, with selection, on the most favourable price in the $m$ blocks before that, until its marginal profit is zero or a pool-wide budget is spent. The budget refills at $q$ of TVL per 12 seconds and holds at most $q$.

\paragraph{C5. Retail flow.} Trades arrive as a Poisson process with intensity (volume/TVL)/(mean size) a year; sizes are lognormal with mean 0.01\% of TVL and log standard deviation 1; the token is drawn in proportion to weight and the side with equal odds. Retail is uninformed.

\paragraph{C6. Stress scenarios.} \emph{Volatility}: all $\sigma_i$ scaled by 0.6 or 1.4. \emph{Regimes}: a two-state Markov chain, the high state with twice the volatility of the low state, 20\% of the time, mean spell 7 days, normalised so average variance is unchanged. \emph{Trends}: a constant drift $\mu_i\sim N(0,0.6^2)$ a year per token and path. \emph{Mean reversion}: an Ornstein--Uhlenbeck log price with $\kappa=2$ a year. \emph{Jumps}: market-wide jumps 26 times a year with standard deviation 4\% scaled by $\sigma_i/\sigma_{\mathrm{ETH}}$, and token-specific jumps 12 times a year per token with standard deviation 6\%, compensated so prices remain martingales.

\paragraph{C7. Estimators.} The direct estimate of loss is the path average of $-\ln(V_T/V_{\mathrm{CRP},T})/T$. For gated designs its noise is dominated by tracking (Table~\ref{tab:gatevalid}), so we also use a control-variate estimate: realised trading costs minus the conditional expected tracking gap $\sum_k[x_k^{\top}\mu-w^{\top}\Sigma x_k-\tfrac12x_k^{\top}\Sigma x_k]\Delta t$ (A12), with $\Sigma$ including jump variance where jumps are on. The two have the same mean; the second removes the martingale part. Standard errors are across paths. The constant $K$ comes from 4{,}000 paths of the Gaussian drift process with covariance $\Omega$ (A11), stepped at $\delta^2/2000K$ years until $\tfrac12\|x\|_1\ge\delta$.

\paragraph{C8. Earlier simulators.} The single-arbitrageur pools and the constant-product comparison of Section~\ref{sec:retail} come from two earlier programs with the same price model, retail model and costs ($c=s=10$\bps, $\delta=2.5\%$, 12-second blocks, 100 paths per scenario).

\paragraph{C9. Reproducibility.} The programs, the commands and seeds for every table and the scripts that turn their output into the tables and figures are available from the author. A full run of all scenarios takes about one hour on two CPU cores.
